\documentclass{zc-ust-hw}

\usepackage{lipsum}

\name{Ali Ashraf,
Abdelrahman Mohamed,
Ahmed Mohamady,
}
\id{202301812,
202300056,
202301826
}
\course{MATH307}
\assignment{\textbf{Fast Fourier Transform (FFT) Implementation and Its Applications.}}

\begin{document}
\maketitle
\tableofcontents
\pagebreak
\section{Abstract}
The Fast Fourier Transform (FFT) is a fundamental numerical algorithm that efficiently converts signals from the time domain to the frequency domain, significantly improving upon the classical Discrete Fourier Transform (DFT), which has a computational complexity of \(O(N^2)\). By using recursive decomposition, the FFT reduces this complexity to \(O(N \log N)\), making it suitable for large-scale problems. This project focuses on implementing and analyzing the FFT, comparing it with the analytical DFT in terms of efficiency and accuracy, and exploring its applications in signal processing, frequency analysis, and image processing. In particular, FFT will be applied to signal filtering and noise reduction by transforming signals to the frequency domain, removing unwanted components, and reconstructing the cleaned signal using the inverse transform, demonstrating its role in communication and audio systems. It will also be used for frequency analysis and feature detection by identifying dominant frequency components in signals, with applications in engineering fields such as vibration analysis, biomedical signal processing, and speech analysis. Additionally, the project emphasizes image processing as an AI-focused application, where the FFT is extended to two-dimensional data to analyze images in the frequency domain; using 2D FFT, images can be filtered to remove noise, enhanced by emphasizing edges and textures, and compressed by discarding less significant frequency components. These techniques form the foundation of AI tasks such as feature extraction, pattern recognition, and preprocessing for deep learning models. Using Python-based simulations, the study evaluates performance under varying conditions such as signal size and noise levels, ultimately demonstrating the FFT's critical role in modern numerical analysis and artificial intelligence applications.


\section*{}

% \begin{thebibliography}{99}
% Empty for now 
% \end{thebibliography}
\end{document}
